% GENERATED FILE. Edit canonical lessons, then run npm run generate:books.
% canonical-source-hash: fnv1a64:7e2a638ff107d67a
% canonical-lessons: JA-C311-omoide, JA-C311-koro, JA-C311-shounen, JA-C311-ushinau, JA-C311-furikaeru, JA-R311-review-307-311

\chapter{Memories of Childhood}
\label{ch:ja-memories-of-childhood}
\hlchaptermodality{311}

\begin{chapteropening}
\textbf{By the end of this chapter:} \emph{I can talk about what I remember from my childhood.}

Complete the last lesson of chapter 311: I can talk about what I remember from my childhood.
\end{chapteropening}

\section[\texorpdfstring{\ja{おもいで} (omoide)}{おもいで (omoide)}]{\ja{おもいで} (omoide) — a memory; a keepsake}

\label{lesson:JA-C311-omoide}



\begin{quote}
Before the new one: say the Japanese for to come true, then the Japanese for shortly.
\end{quote}

\subsection*{You'll want to know: \ja{おもいで}}
\textbf{\ja{おもいで}} — \emph{omoide} — \textquotedblleft{}a memory; a keepsake\textquotedblright{}.

\begin{quote}
\textbf{\ja{なつやすみは} \ja{いい} \ja{おもいでに} \ja{なりました。}} — \emph{natsuyasumi wa ii omoide ni narimashita} — \textquotedblleft{}The summer holidays became a happy memory.\textquotedblright{}
\end{quote}

\begin{cousinweb}
\textbf{\ja{おもいで}} (\emph{omoide}) is a native Japanese word (wago), inherited rather than borrowed from Chinese, built from \emph{omoi}, \textquotedblleft{}thinking\textquotedblright{} (from \emph{omou}), and \emph{de}, \textquotedblleft{}coming out\textquotedblright{} (from \emph{deru}) (JMdict, KANJIDIC2).
\end{cousinweb}

\begin{grammarlens}[title={What you've built}]
One more word for telling what happened.
\end{grammarlens}

\subsection*{Your turn}
\begin{itemize}
\raggedright
  \item \emph{Say it:} \emph{omoide}
  \item \emph{Say it:} \emph{omoide}, once more
  \item \emph{Say it:} \emph{natsuyasumi wa ii omoide ni narimashita} — The summer holidays became a happy memory.
  \item \emph{Recall:} say the Japanese for to come true, then the Japanese for shortly, then say \emph{omoide} again
\end{itemize}

\begin{tcolorbox}[breakable,colback=teal!4,colframe=teal!35!black,title={Before you move on}]
What does \ja{おもいで} mean? (\textquotedblleft{}A memory; a keepsake\textquotedblright{}.) Say it once more.
\end{tcolorbox}
\section[\texorpdfstring{\ja{ころ} (koro)}{ころ (koro)}]{\ja{ころ} (koro) — the time (when); those days}

\label{lesson:JA-C311-koro}



\begin{quote}
Before the new one: say the Japanese for shortly, then the Japanese for a memory.
\end{quote}

\subsection*{You'll want to know: \ja{ころ}}
\textbf{\ja{ころ}} — \emph{koro} — \textquotedblleft{}the time (when); those days\textquotedblright{}.

\begin{quote}
\textbf{\ja{こどもの} \ja{ころ、いなかに} \ja{すんでいました。}} — \emph{kodomo no koro, inaka ni sunde imashita} — \textquotedblleft{}As a child I lived in the countryside.\textquotedblright{}
\end{quote}

\begin{cousinweb}
\textbf{\ja{ころ}} (\emph{koro}) is a native Japanese word (wago), inherited rather than borrowed from Chinese; its kanji was attached later as a kun'yomi (JMdict).
\end{cousinweb}

\begin{grammarlens}[title={What you've built}]
One more word for telling what happened.
\end{grammarlens}

\subsection*{Your turn}
\begin{itemize}
\raggedright
  \item \emph{Say it:} \emph{koro}
  \item \emph{Say it:} \emph{koro}, once more
  \item \emph{Say it:} \emph{kodomo no koro, inaka ni sunde imashita} — As a child I lived in the countryside.
  \item \emph{Recall:} say the Japanese for shortly, then the Japanese for a memory, then say \emph{koro} again
\end{itemize}

\begin{tcolorbox}[breakable,colback=teal!4,colframe=teal!35!black,title={Before you move on}]
What does \ja{ころ} mean? (\textquotedblleft{}The time (when); those days\textquotedblright{}.) Say it once more.
\end{tcolorbox}
\section[\texorpdfstring{\ja{しょうねん} (shōnen)}{しょうねん (shōnen)}]{\ja{しょうねん} (shōnen) — a boy; boyhood}

\label{lesson:JA-C311-shounen}



\begin{quote}
Before the new one: say the Japanese for a memory, then the Japanese for the time (when).
\end{quote}

\subsection*{You'll want to know: \ja{しょうねん}}
\textbf{\ja{しょうねん}} — \emph{shōnen} — \textquotedblleft{}a boy; boyhood\textquotedblright{}.

\begin{quote}
\textbf{\ja{しょうねんの} \ja{ころ、そふは} \ja{やまで} \ja{はたらいていました。}} — \emph{shōnen no koro, sofu wa yama de hataraite imashita} — \textquotedblleft{}As a boy, my grandfather worked in the mountains.\textquotedblright{}
\end{quote}

\begin{cousinweb}
\textbf{\ja{しょうねん}} (\emph{shōnen}) is Sino-Japanese (kango): two Chinese characters read with their on'yomi, readings borrowed from Chinese. They are \emph{shō}, \textquotedblleft{}little\textquotedblright{} (Middle Chinese \emph{syewH}) and \emph{nen}, \textquotedblleft{}year\textquotedblright{} (\emph{nen}) (KANJIDIC2; Middle Chinese in Baxter's transcription).
\end{cousinweb}

\begin{grammarlens}[title={What you've built}]
One more word for telling what happened.
\end{grammarlens}

\subsection*{Your turn}
\begin{itemize}
\raggedright
  \item \emph{Say it:} \emph{shōnen}
  \item \emph{Say it:} \emph{shōnen}, once more
  \item \emph{Say it:} \emph{shōnen no koro, sofu wa yama de hataraite imashita} — As a boy, my grandfather worked in the mountains.
  \item \emph{Recall:} say the Japanese for a memory, then the Japanese for the time (when), then say \emph{shōnen} again
\end{itemize}

\begin{tcolorbox}[breakable,colback=teal!4,colframe=teal!35!black,title={Before you move on}]
What does \ja{しょうねん} mean? (\textquotedblleft{}A boy; boyhood\textquotedblright{}.) Say it once more.
\end{tcolorbox}
\section[\texorpdfstring{\ja{うしなう} (ushinau)}{うしなう (ushinau)}]{\ja{うしなう} (ushinau) — to lose for good; to be robbed of}

\label{lesson:JA-C311-ushinau}



\begin{quote}
Before the new one: say the Japanese for the time (when), then the Japanese for a boy.
\end{quote}

\subsection*{You'll want to know: \ja{うしなう}}
\textbf{\ja{うしなう}} — \emph{ushinau} — \textquotedblleft{}to lose for good; to be robbed of\textquotedblright{}.

\begin{quote}
\textbf{\ja{ちちは} \ja{じしんで} \ja{しごとを} \ja{うしないました。}} — \emph{chichi wa jishin de shigoto o ushinaimashita} — \textquotedblleft{}My father lost his job in the earthquake.\textquotedblright{}
\end{quote}

Polite \textbf{\ja{うしないます}}, \emph{ushinaimasu}; a godan verb. Past (polite): \textbf{\ja{おとこが} \ja{しごとを} \ja{うしないました。}}, \emph{otoko ga shigoto o ushinaimashita}.

\begin{cousinweb}
\textbf{\ja{うしなう}} (\emph{ushinau}) is a native Japanese word (wago), inherited rather than borrowed from Chinese; its kanji was attached later as a kun'yomi (JMdict).
\end{cousinweb}

\begin{grammarlens}[title={What you've built}]
One more word for telling what happened.
\end{grammarlens}

\subsection*{Your turn}
\begin{itemize}
\raggedright
  \item \emph{Say it:} \emph{ushinau}
  \item \emph{Say it:} \emph{ushinau}, once more
  \item \emph{Say it:} \emph{chichi wa jishin de shigoto o ushinaimashita} — My father lost his job in the earthquake.
  \item \emph{Recall:} say the Japanese for the time (when), then the Japanese for a boy, then say \emph{ushinau} again
\end{itemize}

\begin{tcolorbox}[breakable,colback=teal!4,colframe=teal!35!black,title={Before you move on}]
What does \ja{うしなう} mean? (\textquotedblleft{}To lose for good; to be robbed of\textquotedblright{}.) Say it once more.
\end{tcolorbox}
\section[\texorpdfstring{\ja{ふりかえる} (furikaeru)}{ふりかえる (furikaeru)}]{\ja{ふりかえる} (furikaeru) — to look back (on)}

\label{lesson:JA-C311-furikaeru}



\begin{quote}
Before the new one: say the Japanese for a boy, then the Japanese for to lose for good.
\end{quote}

\subsection*{You'll want to know: \ja{ふりかえる}}
\textbf{\ja{ふりかえる}} — \emph{furikaeru} — \textquotedblleft{}to look back (on)\textquotedblright{}.

\begin{quote}
\textbf{\ja{そふは} \ja{ときどき} \ja{むかしを} \ja{ふりかえります。}} — \emph{sofu wa tokidoki mukashi o furikaerimasu} — \textquotedblleft{}Grandfather sometimes thinks back to the old days.\textquotedblright{}
\end{quote}

Polite \textbf{\ja{ふりかえります}}, \emph{furikaerimasu}; a godan verb. Past (polite): \textbf{\ja{ははが} \ja{ふりかえりました。}}, \emph{haha ga furikaerimashita}.

\begin{cousinweb}
\textbf{\ja{ふりかえる}} (\emph{furikaeru}) is a native Japanese word (wago), inherited rather than borrowed from Chinese, built from \emph{furi}, \textquotedblleft{}swinging\textquotedblright{} (from \emph{furu}), and \emph{kaeru}, \textquotedblleft{}to turn back\textquotedblright{}, the intransitive partner of \emph{kaesu}, \textquotedblleft{}to give back\textquotedblright{} (JMdict, KANJIDIC2).
\end{cousinweb}

\begin{grammarlens}[title={What you've built}]
One more word for telling what happened.
\end{grammarlens}

\subsection*{Your turn}
\begin{itemize}
\raggedright
  \item \emph{Say it:} \emph{furikaeru}
  \item \emph{Say it:} \emph{furikaeru}, once more
  \item \emph{Say it:} \emph{sofu wa tokidoki mukashi o furikaerimasu} — Grandfather sometimes thinks back to the old days.
  \item \emph{Recall:} say the Japanese for a boy, then the Japanese for to lose for good, then say \emph{furikaeru} again
\end{itemize}

\begin{tcolorbox}[breakable,colback=teal!4,colframe=teal!35!black,title={Before you move on}]
What does \ja{ふりかえる} mean? (\textquotedblleft{}To look back (on)\textquotedblright{}.) Say it once more.
\end{tcolorbox}
\section[(dialogue)]{Looking back: from A Call in the Night to Memories of Childhood}

\label{lesson:JA-R311-review-307-311}



\begin{quote}
Say the Japanese for to be connected and to join.
\end{quote}

\subsection*{Your turn}
Say these aloud:

\begin{itemize}
\raggedright
  \item to be connected, to join, a message (passed on), a mobile phone, to call back
  \item a mistake, to be heard, a telephone receiver, to speak to (someone), busy
  \item to give to me, to give (to someone else), a gift in return, a person, to confide
  \item all along, actually, in the end, to come true, shortly
  \item a memory, the time (when), a boy, to lose for good, to look back (on)
\end{itemize}

\begin{tcolorbox}[breakable,colback=teal!4,colframe=teal!35!black,title={Before you move on}]
To be connected? (\textbf{\ja{つながる}}, \emph{tsunagaru}.) To join? (\textbf{\ja{つなぐ}}, \emph{tsunagu}.) A message (passed on)? (\textbf{\ja{でんごん}}, \emph{dengon}.) A mobile phone? (\textbf{\ja{けいたい}}, \emph{keitai}.) To call back? (\textbf{\ja{かけなおす}}, \emph{kakenaosu}.) A mistake? (\textbf{\ja{まちがい}}, \emph{machigai}.) To be heard? (\textbf{\ja{きこえる}}, \emph{kikoeru}.) A telephone receiver? (\textbf{\ja{じゅわき}}, \emph{juwaki}.) To speak to (someone)? (\textbf{\ja{はなしかける}}, \emph{hanashikakeru}.) Busy? (\textbf{\ja{いそがしい}}, \emph{isogashii}.) To give to me? (\textbf{\ja{くれる}}, \emph{kureru}.) To give (to someone else)? (\textbf{\ja{あげる}}, \emph{ageru}.) A gift in return? (\textbf{\ja{おかえし}}, \emph{okaeshi}.) A person? (\textbf{\ja{ひと}}, \emph{hito}.) To confide? (\textbf{\ja{うちあける}}, \emph{uchiakeru}.) All along? (\textbf{\ja{ずっと}}, \emph{zutto}.) Actually? (\textbf{\ja{じつは}}, \emph{jitsu wa}.) In the end? (\textbf{\ja{とうとう}}, \emph{tōtō}.) To come true? (\textbf{\ja{かなう}}, \emph{kanau}.) Shortly? (\textbf{\ja{まもなく}}, \emph{mamonaku}.) A memory? (\textbf{\ja{おもいで}}, \emph{omoide}.) The time (when)? (\textbf{\ja{ころ}}, \emph{koro}.) A boy? (\textbf{\ja{しょうねん}}, \emph{shōnen}.) To lose for good? (\textbf{\ja{うしなう}}, \emph{ushinau}.) To look back (on)? (\textbf{\ja{ふりかえる}}, \emph{furikaeru}.)
\end{tcolorbox}
